\emph{Les trois questions sont indépendantes.}

\subsection*{Question 1}

	On considère l'intégrale suivante :
	\[I = \displaystyle\int_{-2}^0 \left(5\e^{-2t} + 3\right)\:\text{d}t.\]

	Soit $f$ la fonction définie sur $\R$ par $f(t) = 5\e^{-2t} + 3$.

	Soit $F$ la fonction définie sur $\R$ par $F(t) = - 2,5\e^{-2t} + 3t$.

	On admet que $F$ est dérivable sur $\R$.


	\begin{enumerate}
		\item Démontrer que la fonction $F$ est une primitive de la fonction $f$ sur $\R$.

		\item Démontrer que $I = 3,5 + 2,5\e^4$.
	\end{enumerate}

\subsection*{Question 2}
	On considère la fonction $g$ définie sur $\R$ par $g(x) = 2\e^{3x} - 16$.

	On admet que $g$ est dérivable sur $\R$ et on note $g'$ sa fonction dérivée.

	\begin{enumerate}
		\item Démontrer que pour tout réel $x$,\quad $g'(x) > 0$.

		\item En déduire le sens de variation de la fonction $g$ sur $\R$.

		\item Résoudre, en détaillant les étapes, l'équation $g(x) = 0$.
	\end{enumerate}

\subsection*{Question 3}
	On désigne par i le nombre complexe de module 1 et d'argument $\pi$.

	Le plan complexe est rapporté à un repère orthonormé \Ouv.

	On considère les points A et B d'affixes respectives
	$z_{\text{A}} = 4\e^{\text{i}\frac{5\pi}{3}}$ et
	$z_{\text{B}} = 2\e^{\text{i}\frac{2\pi}{3}}$.

	On rappelle l'écriture exponentielle suivante : $-2 = 2\e^{\text{i}\pi}$.

\medskip

	\begin{enumerate}
		\item Démontrer, en détaillant les calculs, que $z_{\text{A}} = - 2 \times z_{\text{B}}$.
		\item Parmi les quatre propositions suivantes, une seule est vraie.

		Indiquer sur la copie le numéro de celle-ci. Aucune justification n'est demandée.

		\medskip

		\textbf{Proposition 1 :} A est l'image de B par la rotation de centre O et d'angle $- \pi$.

		\textbf{Proposition 2 :} A est l'image de B par l'homothétie de centre O et de rapport $-2$.

		\textbf{Proposition 3 :} A est l'image de B par la translation de vecteur $\vect{u}$ d'affixe $-2$.

		\textbf{Proposition 4 :} A est l'image de B par l'homothétie de centre O et de rapport 2.

	\end{enumerate}
